\section{Proofs and Algebra}\label{app:proofs}

\subsection{Global continuation from primitives}\label{app:global-continuation}

Write $A=a_L$, $d=a_H-a_L>0$, and $n=l+h$. For a fixed marginal price $p$, provider profit is affine and strictly increasing in the fixed fee $F$ whenever a positive mass continues to participate. Hence, within any fixed participation set, an optimal finite $F$ is the largest fee that preserves that set. Negative fixed fees cannot be optimal for a positive served mass: raising such an $F$ until a participation constraint binds weakly preserves the active set and strictly raises profit. Weak participation is used here to attain a binding threshold exactly.

Because $S_H(p)\geq S_L(p)$ for every $p$, an anonymous tariff cannot serve an installed $L$ while excluding an installed $H$. When both masses are positive, the only economically distinct active sets are therefore both served, $H$ only, and no service. A zero-$h$ state has no separate screening issue: ``$L$ only'' then simply means serving the only installed type. Strict-$\mathcal R^+$ equilibrium candidates satisfy $h\geq h_0>0$.

For $0\leq p<a_L$, both types have positive usage. The highest fixed fee serving both is $F=S_L(p)$, so
\[
 B(p)\equiv\Pi_M^B(p\mid l,h)
 =n\frac{(A-p)^2}{2}+(p-c)\{l(A-p)+h(A+d-p)\}.
\]
The candidate price is
\[
 p^*=c+\frac{dh}{n}.
\]
Direct completion of the square, rather than a local FOC, gives the exact identity
\begin{equation}
 B(p^*)-B(p)=\frac{n}{2}(p-p^*)^2\geq0.
 \label{eq:square-loss}
\end{equation}
Thus $p^*$ is the unique global maximizer on any both-served price interval containing it. Under strict $\mathcal R^+$, $p^*>c>0$ and $p^*(h)\leq p^*(h_F)<a_L$, so the metered optimum is attained despite the open architecture restriction $p>0$.

For an $H$-only outcome at a price $p<a_H$, the optimal fee is $F=S_H(p)$ and per-$H$ profit is
\[
 S_H(p)+(p-c)(a_H-p).
\]
This equals total surplus net of inference cost and is uniquely maximized at $p=c$, yielding
\[
 \Pi_M^H(h)=h\frac{(a_H-c)^2}{2}.
\]
Since $0<c<a_L$, this maximum is feasible and attained in the metered architecture. Under flat pricing, $p=0$ and the analogous $H$-only value is equation~\eqref{eq:flat-H-profit}.

The price boundaries introduce no missing branch. At $p=a_L$, type $L$ has zero usage and $S_L=0$; if $F=0$ it participates by the weak-IR convention but contributes neither usage nor a positive fixed payment, so the economic allocation is the $H$-only boundary. For $a_L<p<a_H$, any positive $F$ excludes $L$ and the problem is $H$ only; choosing $F\leq0$ merely adds a zero-use $L$ mass paying a nonpositive fee and cannot improve on the corresponding $H$-only tariff. For $p\geq a_H$, both usage quantities are zero, and no service weakly dominates any subsidy. These observations also cover arbitrary finite negative-$F$ deviations.

It remains to compare the globally optimized both-served and $H$-only branches. At the metered optimum, direct simplification gives
\begin{equation}
 D_M(h)\equiv B(p^*(h))-\Pi_M^H(h)
 =\frac{l(A-c)^2}{2}-d(A-c)h-\frac{d^2lh}{2(l+h)},
 \label{eq:DM}
\end{equation}
with
\begin{equation}
 D'_M(h)=-d(A-c)-\frac{d^2l^2}{2(l+h)^2}<0.
 \label{eq:DMprime}
\end{equation}
For flat pricing,
\begin{equation}
 D_F(h)\equiv\Pi_F^B(l,h)-\Pi_F^H(h)
 =l\left(\frac{A^2}{2}-cA\right)-\frac{hd(2A+d)}{2},
 \label{eq:DF}
\end{equation}
so
\begin{equation}
 D'_F(h)=-\frac{d(2A+d)}{2}<0.
 \label{eq:DFprime}
\end{equation}
Strict conditions \eqref{eq:Rplus6}--\eqref{eq:Rplus7} say $D_F(h_F)>0$ and $D_M(h_F)>0$. Monotonicity therefore gives strict both-served dominance over $H$ only at every $h\leq h_F$. No service is also dominated throughout the candidate interval: $\Pi_F^B$ is affine in $h$ and is strictly positive at both endpoints by \eqref{eq:Rplus5}, while optimized metered both-served profit equals flat both-served profit plus the strictly positive gross gain $\Phi(l,h)$. Hence equations \eqref{eq:VF-global}--\eqref{eq:VM-global} select the both-served branches throughout the complete candidate interval.

\subsection{Why the candidate state interval is complete}\label{app:candidate-completeness}

The preceding global problem implies the state bounds rather than assuming them. At every optimal continuation, type $L$ receives zero continuation rent. If both types are served, $F=S_L(p)$ binds $L$'s participation constraint. If only $H$ is served or there is no service, $L$ receives zero. Thus its integration condition is $k\leq b_L$, and the uniform cost distribution plus $l\in(0,1)$ gives $l=b_L/K_L$ in every candidate equilibrium.

Type $H$'s optimal continuation rent lies in $[0,R_F]$. An $H$-only optimum extracts all of its continuation surplus, and no service gives zero. On a both-served branch with $0\leq p\leq a_L$,
\[
 S_H(p)-S_L(p)=R_F-dp\in[0,R_F],
\]
where flat pricing $p=0$ attains the upper endpoint. Any provider mixture over optimal continuations yields an expected rent in the same closed interval. Since strict $\mathcal R^+$ imposes $b_H+R_F<K_H$, the uniform-cost cutoff never saturates, and therefore
\[
 h=\frac{b_H+\mathbb E[R_H]}{K_H}\in
 \left[\frac{b_H}{K_H},\frac{b_H+R_F}{K_H}\right]
 =[h_0,h_F].
\]
This establishes candidate-interval completeness before the endpoint-dominance argument of the previous subsection is invoked. Positive $b_H$ gives $h_0>0$. Exact fixed-fee attainment uses weak participation; under strict participation an $\varepsilon$-implementation or limit argument would be a different model and is not claimed here.

\subsection{Proof of Proposition~\ref{prop:T1}}

Define
\[
g(h)=b_H+R_F-d\left(c+d\frac{h}{l+h}\right)-K_Hh.
\]
The metered rational-expectations state is the zero of $g$. Since $K_Hh_0=b_H$,
\[
g(h_0)=R_F-dp^*(h_0).
\]
The function $p^*(h)$ is increasing in $h$, and strict $\mathcal R^+$ gives $p^*(h_F)<a_L$. Hence $p^*(h_0)<a_L$. Moreover,
\[
\frac{R_F}{d}=\frac{a_H+a_L}{2}>a_L,
\]
so $p^*(h_0)<R_F/d$ and therefore $g(h_0)>0$. Using $K_Hh_F=b_H+R_F$,
\[
g(h_F)=-d p^*(h_F)<0.
\]
Finally,
\[
g'(h)=-K_H-\frac{d^2l}{(l+h)^2}<0.
\]
Thus $g$ has exactly one zero and it lies in $(h_0,h_F)$. Hence $h_0<h_M<h_F$. \qed

\subsection{Proof of Proposition~\ref{prop:T2}}

Equation \eqref{eq:Phi-derivative} implies that $\Phi(l,h)$ is strictly increasing in $h$. Proposition~\ref{prop:T1} gives $h_M<h_F$. Therefore
\[
\mu_M=\Phi(l,h_M)<\Phi(l,h_F)=\mu_F.
\]
\qed

\subsection{Pure boundaries and proof of Proposition~\ref{prop:T3}}\label{app:mixed-proof}

Under a pure metered expectation the installed state is $h_M$. The provider weakly prefers metering exactly when $\mu\leq\mu_M$, so the pure metered outcome is self-consistent on that closed region. Under a pure flat expectation the state is $h_F$, and flat pricing is weakly optimal exactly when $\mu\geq\mu_F$.

The equalities require an aggregation check. At $\mu=\mu_M$, suppose the provider used a metering probability $\rho<1$. The integration equation evaluated at $h_M$ has excess right-hand side
\[
 (1-\rho)d p^*(h_M)>0.
\]
Its left-minus-right crossing is strictly increasing because $K_H+\rho d\,p^{*\prime}(h)>0$, so the induced root satisfies $h>h_M$. Strict monotonicity of $\Phi$ then gives $\Phi(l,h)>\mu_M$, making metering strictly optimal, a contradiction to $\rho<1$. Hence the aggregate boundary equilibrium has $\rho=1$ and $h=h_M$.

At $\mu=\mu_F$, any $\rho>0$ implies from
\[
 K_H(h_F-h)=\rho d p^*(h)>0
\]
that $h<h_F$. Then $\Phi(l,h)<\mu_F$, so flat pricing is strictly optimal and positive metering probability is impossible. Thus the aggregate boundary equilibrium has $\rho=0$ and $h=h_F$. Provider indifference at either pure induced state therefore does not generate a continuum of aggregate mixtures.

For the strict interval $\mu_M<\mu<\mu_F$, neither pure expectation is self-consistent. Because $\Phi(l,h)$ is continuous and strictly increasing, there is a unique $h^*\in(h_M,h_F)$ satisfying $\Phi(l,h^*)=\mu$. At that state the provider is indifferent between the globally optimal flat tariff
\[
 (F_F,p_F)=\left(\frac{a_L^2}{2},0\right)
\]
and the globally optimal metered tariff
\[
 (F_M,p_M)=\left(S_L(p^*(h^*)),p^*(h^*)\right).
\]
If the provider realizes the metered tariff with probability $\rho$, the high-use integration condition is
\[
K_Hh^*=b_H+R_F-\rho d p^*(h^*),
\]
so
\[
\rho^*=\frac{K_H(h_F-h^*)}{d p^*(h^*)}.
\]
Positivity follows from $h^*<h_F$. Since $g$ is strictly decreasing and $h^*>h_M$, $g(h^*)<0$, which is equivalent to
\[
K_H(h_F-h^*)<d p^*(h^*).
\]
Thus $0<\rho^*<1$.

The user strategies in \eqref{eq:integration-strategies} are best responses because $L$ obtains zero continuation rent under either realized tariff, while $H$ obtains $R_F$ under flat pricing and $R_F-dp^*(h^*)$ under metering. Hence the expected integration gains are exactly the two cutoff expressions. Uniform costs aggregate these strategies to $l=b_L/K_L$ and $h=h^*$. Users choose future quantities only after the tariff realization; no average-price demand is used. Because users are atomless, a single integration deviation does not change the aggregate state and hence does not change the provider's continuation. Off path, the provider selects the full global continuation in equations \eqref{eq:VF-global}--\eqref{eq:VM-global}, with arbitrary mixing only at genuine payoff ties. Therefore $(h^*,\rho^*)$ is the unique aggregate mixed resolution, while measure-zero cutoff actions and off-path tie-breaking need not be unique. \qed

\subsection{Proof of Corollary~\ref{cor:collapse}}

If architecture does not change the installed high-use state, write that common state as $\bar h$. Then
\[
\mu_M=\Phi(l,\bar h)=\mu_F,
\]
so the strict interval between the thresholds is empty. \qed

\subsection{Proof of Proposition~\ref{prop:W1}}

At a fixed installed state,
\[
\mu_P=\frac{l+h}{2}\left(c+d\frac{h}{l+h}\right)^2.
\]
Using equation \eqref{eq:omega}, the welfare gain from moving from flat pricing to the provider-optimal metered continuation, before paying $\mu$, is
\[
\mu_W=\frac{l+h}{2}\left[c^2-d^2\left(\frac{h}{l+h}\right)^2\right].
\]
Subtracting yields
\[
\mu_P-\mu_W
=dh\left(c+d\frac{h}{l+h}\right)
=dh\,p^*(h)>0.
\]
\qed

\subsection{Exact outside-$\mathcal R^+$ scope guard}

For the parameter vector in \eqref{eq:scope-counterexample}, the both-served flat continuation sets $F=a_L^2/2=8$ and earns
\[
(0.1+0.6)8-[0.1(4)+0.6(5)]=\frac{11}{5}.
\]
An $H$-only flat continuation sets $F=a_H^2/2=25/2$ and earns
\[
0.6\left(\frac{25}{2}-5\right)=\frac{9}{2}.
\]
Hence the $H$-only advantage is
\[
\frac{9}{2}-\frac{11}{5}=\frac{23}{10}>0.
\]
This is the permanent 23/10 scope guard. The calculation establishes only that the both-served candidate need not be globally optimal outside the certified domain; it does not characterize the full equilibrium correspondence there.
